\chapter{Conclusions and Future Work}\label{chapter-06}
This dissertation addressed the challenge of introducing reconfigurability into static structural systems without abandoning the load-bearing function that makes these systems useful. The central premise of the work is that reconfigurability in structural applications cannot be treated solely as a mechanism-design problem; reconfigurability must also be coupled with strength, stability, fabrication constraints, and predictable load transfer. This dissertation considered the inverse problem of starting with established structural forms and systematically introducing mobility, deployability, and programmable mechanical response. This philosophy was pursued through four objectives: developing a framework for transforming static trusses into reconfigurable systems using planar quadrilateral linkage principles; developing computational tools for evaluating the kinematic and mechanical response of the transformed systems; incorporating programmed hinge behavior to enable functional responses such as multistability and reusable energy-management structures; and broadening the framework beyond civil truss systems through curved-crease origami hulls that are offer rapid deployability and tunable hydrodynamic performance. Together, these objectives establish a foundation for the design and analysis of reconfigurable structural systems that can change shape while retaining meaningful structural function across multiple configurations.

\section{Key Contributions of the Dissertation}\label{06-sec-contributions}
The key contributions of this dissertation are summarized below.

\noindent{}\textbf{A framework for transforming static trusses into reconfigurable systems and evaluating their kinematic behavior.}
This dissertation developed a systematic framework for converting static trusses into reconfigurable systems using principles of flat-foldable quadrilateral linkages. The method transforms triangular units of static trusses into quadrilateral linkages by introducing an additional node into selected members. The node location is determined using closed-form expressions, derived for all triangle geometries, that ensure the resulting quadrilateral linkage satisfies the Grashof flat-foldability criterion (Fig.~\ref{fig-node-intro}). These nodes are placed on tensile members so that the transformed truss retains a stable load path and can carry the prescribed design load despite having mobility greater than one. The unit-level node-placement rules are extended to a system-level workflow that selects candidate members based on both the geometry of the triangular units and the corresponding member forces. The workflow applies to most statically determinate trusses with a continuous sequence of triangular units between the top and bottom chords. Its applicability was demonstrated using conventional truss layouts, including Scissor, Fink, and Warren trusses, as well as topology-optimized layouts, including cantilever, L-shaped, and serpentine trusses. A sequential kinematic simulator was also developed to visualize deployment paths and quantify packing performance as the transformed trusses morph through sequential actuation of their kinematic degrees of freedom. The resulting reconfigurable trusses can reduce their convex hull area by up to 93\% and achieve compact storage while preserving the structural role of the original truss topology.

\noindent{}\textbf{Proof-of-concept prototypes to demonstrate structural capacities of reconfigurable trusses.} We validated the feasibility of constructing real-life reconfigurable trusses using our method by fabricating two proof-of-concept prototypes \textemdash{} a 30~cm reconfigurable cantilever truss and a 3~m Warren truss. Both prototypes employ a joint design strategy that distributes truss members in and out of the plane to ensure desired shape-morphing behavior without member contact or overlap. Load testing of the cantilever trusses confirmed that the reconfigurable truss exhibits force-stable behavior and maintains a peak load capacity and structural stiffness comparable to its static counterparts despite having mobility greater than one. The three-meter-long reconfigurable Warren truss weighing 38~kg supported a peak load of 19~kN, demonstrating the structural capacity and scalability of the proposed designs.

\noindent{}\textbf{Computational tools for mechanical analysis of reconfigurable trusses.}
This dissertation developed computational tools for efficient nonlinear analysis of reconfigurable bar-linked structures whose global response is governed by joint rotational stiffness. The Three-Node Torsional Spring (3NTS) element provides a compact way to incorporate torsional springs into truss-like models without introducing rotational degrees of freedom or replacing joints with artificial beam elements. As a result, reconfigurable trusses with torsional springs can be analyzed using the same translational degrees of freedom as conventional truss models, while capturing large-displacement behavior, elastic buckling, and postbuckling response. This formulation provides the mechanical analysis capability needed to move beyond purely kinematic studies and evaluate how hinge stiffness influences force response, stability, and functional reconfiguration.

\noindent{}\textbf{Programmable hinge mechanics for functional reconfigurable trusses.}
This dissertation extended the modeling capabilities of mechanical analysis tools to enable programmable mechanical response in reconfigurable trusses. An enhanced torsional spring law was introduced to bound the admissible rotation of the 3NTS element while preserving its intended linear response range. This law enables the simulation of sequential folding, locking, and load redistribution among the kinematic degrees of freedom of the truss. The framework was further expanded through magnetic multistable and idealized bistable hinge models, which introduce repeated energy barriers, multiple stable rotational states, and snap-through transitions in the reconfigurable trusses. These models show that local hinge behavior can be programmed to achieve a desired global structural response. In particular, reconfigurable trusses can manage impact energy through controlled geometric transformations rather than permanent material damage. This allows the primary members to remain elastic while the structure absorbs external work through recoverable reconfiguration. This capability provides a basis for reusable impact-energy management systems, controlled deployment paths, and intermediate stable configurations relevant to robotic structures, deployable space systems, and support frames for membrane or origami-based systems.

\noindent{}\textbf{Rapidly deployable, shape-morphing curved-origami hulls with tunable hydrodynamic performance.} Curved-crease origami hulls introduced in this dissertation can reproduce the hydrodynamic behavior of conventional planing hulls while also enabling geometry-based performance tuning. Because these hulls can closely conform to conventional hull geometries, their hydrodynamic response closely matches that of their conventional counterparts, including the VPS and GPPH hulls. Hydrodynamic simulations performed in Powersea show nearly equivalent trends in friction drag, total resistance, hull sinkage, and trim over a wide range of forward speeds in still water. In head-sea conditions, the curved-crease origami hulls also show close agreement with conventional hulls in the amplitude and frequency of hull motion time histories. These results indicate that hulls designed using curved-origami principles can preserve important performance characteristics associated with power efficiency, passenger comfort, and operational safety while providing the fabrication, storage, and deployment advantages of origami-based systems. In addition to replicating conventional hull behavior, the curved-crease origami hulls enable on-demand tuning of hydrodynamic performance by controlling the actuation of fold lines along the keel and chine. Changing the folding angle modifies key geometric features of the hull, especially the deadrise angle. Lower-deadrise configurations produce flatter planing surfaces that are better suited for calm water because they generally reduce sinkage, trim, and resistance. However, these configurations can produce larger heave and pitch in waves and are therefore less effective for seakeeping in wavy conditions. Higher-deadrise configurations produce sharper V-shaped planing surfaces that generally increase resistance, sinkage, and trim in still water but reduce heave and pitch amplitudes in waves. A single curved-crease origami hull can therefore transition between different deadrise configurations suited for different sea states, allowing the same structure to balance operational efficiency, passenger comfort, and seakeeping performance through controlled shape morphing.

\section{Considerations for Future Work}\label{06-sec-future-work}
While the reconfigurable trusses and origami-hulls introduced in this work offer exciting avenues for designing shape-morphing structures, several challenges must be overcome before translating these designs into practical structures. Future work can focus on tackling the following challenges:

\noindent{}\textbf{General methods for designing reconfigurable trusses for three-dimensional structures.}
The node-introduction workflow described in Section~\ref{02-subsec-algo} provides an effective method for transforming many two-dimensional statically determinate trusses into reconfigurable systems. However, the current implementation relies on prescribed rules for hinge placement and is best suited to trusses formed by a single sequence of triangular units. These assumptions limit its use for trusses with multiple layers of triangles, repeated triangular patterns, arbitrary polygonal cells, or three-dimensional connectivity. In such cases, introducing a new node can alter the mobility of several adjacent units. An unsuitable node or hinge location can therefore produce incompatible kinematics, reduce packing efficiency, or compromise structural stability. Future work should develop a more general design framework that identifies node locations, bar connectivity, and hinge-placement strategies without relying on case-specific rules. Graph-theoretic methods offer a promising direction for this generalization because trusses can be represented through nodes, edges, and connectivity constraints. Such a framework could identify the structural subgraphs requiring mobility, determine which bars should remain continuous for load transfer, and select hinge locations that preserve compatible deployment. This approach could also support the design of reconfigurable trusses generated by topology optimization, whose final layouts often include irregular connectivity and polygons with more than three sides. More broadly, a graph-based formulation would help extend the present framework from planar trusses to complex two-dimensional and three-dimensional truss systems while maintaining a direct connection between kinematics, load paths, and structural performance.

\noindent{}\textbf{Mitigating structural overlap to optimize packing density.} As the complexity and degrees of freedom for reconfigurable trusses increase, the potential for structural self-overlapping during actuation becomes increasingly evident. This issue is particularly highlighted in the Scissor truss example illustrated in Fig.~\ref{fig-trad-truss}\,(a), where the actuation of the third degree of freedom results in portions of the structure overlapping. Real-life truss structures cannot achieve a compact state without providing special provisions for folding mechanisms. For such structures, it is crucial to understand their spatial movement patterns and identify the degrees of freedom whose actuation triggers a self-overlap. Addressing this challenge necessitates future investigation into varied sequences of degree-of-freedom actuation to mitigate structural overlap. Furthermore, automating the distribution of truss members out of the plane and incorporating additional methods for node placement (compressive members) will enhance the feasibility of reconfigurable structures. These practices can be explored to prevent self-overlapping and optimize the ability of the structure to minimize its footprint after shape transformations.


\noindent{}\textbf{Coupled hydrodynamic and structural evaluation of shape-morphing curved-crease origami hulls}. Future work should evaluate the coupled relationship between shape morphing, hydrodynamic loading, and structural response in curved-crease origami hulls. The shape change of these curved-origami hulls alters the deadrise angle and other geometric and inertial properties, including hull length, bow curvature, beam width, center of gravity, and radius of gyration. These changes can alter the hull surface pressure distribution, wave response, and local impact loads experienced by the hull. This effect is especially important near the bow, where planing hulls can experience high local pressures under wave-impact loading. The structural response must also be assessed because the shell thickness may be limited by the curvature and foldability requirements of the origami geometry, as shown in Fig.~\ref{fig-prototype}\,(a). Future studies should therefore examine how payload, deadrise configuration, and hull geometry influence hydrodynamic pressures, global forces, stresses, and strains. Higher-fidelity computational models can be used for this purpose, including computational fluid dynamics and one-way coupled fluid-structure interaction analyses. These models can evaluate quantities that are not fully captured by simplified planning-hull tools, such as added wave resistance, spray resistance, local accelerations, and pressure variations over the hull surface. The relationship between shape morphing and seakeeping performance can also be studied using response amplitude operators, such as heave and pitch transfer functions, across different wavelengths and hull configurations. These analyses would clarify how changes in the curved-crease origami shape affect resonant behavior, impact loading, and structural demand. Two-tank experiments can further support the computational results by providing pressure, force, motion, and acceleration measurements for selected configurations.